\include{beamer}
\documentclass[xcolor=dvipsnames]{beamer} 

\usepackage{listings}
\usecolortheme[named=Green]{structure} 
\usetheme[height=10mm]{Rochester}
\setbeamertemplate{navigation symbols}{}
\setbeamertemplate{footline}{
    \hfill \includegraphics[scale=.20]{Emory_Logo_Blue}\\
    \vspace{1mm}
} 

\begin{document}
\begin{frame}
\begin{center}
\vspace{5mm}
{\huge An Introduction to Git\\ \vspace{1mm} Version Control}\\
\vspace{7mm}
\includegraphics[scale=.3]{gitlogo}\\
\vspace{8mm}
{\large February 4, 2014}\\
Clarissa Garvey, Emory University
\end{center}
\end{frame}

\setbeamertemplate{footline}{
    \hspace{1mm} \includegraphics[scale=.25]{gitlogo}\hfill \includegraphics[scale=.20]{Emory_Logo_Blue}\\
    \vspace{1mm}
} 
\begin{frame}{Overview}

{\LARGE Contents:}
\begin{enumerate}
\item What is Git?
\item Why you should use Git %and where to get more info
\item How Git works
\item Getting Git
\item Basic Git commands
\item Popular Git hosting websites
\item Questions
\end{enumerate}
\end{frame}

\begin{frame}{What is Git?}

\vspace{-5mm}
\begin{itemize}
\item Version Control Systems (VCSs)\\
\hspace{4mm} ``Using a VCS (...) means that if you screw things up or\\\hspace{6mm}  lose files, you can generally recover easily.''
\begin{itemize}
\item Backup
\item Revert state
\item Collaborate
\item Low overhead
\end{itemize}

\item Git is a Distributed VCS
\begin{itemize}
\item Each user has a local repository
\end{itemize}
\end{itemize}

\end{frame}

\begin{frame}{Why you should use Git}
\begin{itemize}
\item Easy to quickly change back to old versions of a file
\item Easy to see differences between file versions
\item Enables experimentation on source code
\item Easy collaboration, even in the same file
\item Allows advisers to easily see progress in projects
\item Easy to share work with others
\item Allows for quick annotation of changes to files
\item It's {\LARGE great!}
\end{itemize}
\end{frame}

\begin{frame}{How Git works}
\vspace{-10mm}
{\LARGE How Git Handles Data}

\begin{itemize}
\item Data as snapshots
\begin{itemize}
\item Contrasts with $\Delta$s
\end{itemize}
\item Files are stored locally
\begin{itemize}
\item Current files in base directory
\item Metadata in .git file
\end{itemize}
\item Almost everything is undoable
\end{itemize}

\end{frame}

\begin{frame}{How Git works}

{\LARGE 3-Tier Model}\\
\begin{itemize}
\item Tracked vs. untracked files

\item {\large Modified}
\begin{itemize}
\item There are unstaged changes to files since the last commit
\end{itemize}

\item {\large Staged}
\begin{itemize}
\item A snapshot of the file has been taken with the intent to commit later
\end{itemize}
\item {\large Committed}
\begin{itemize}
\item The changes have been saved permanently in the .git file with added metadata
\end{itemize}
\end{itemize}

\end{frame}

\begin{frame}{Getting Git}

{\LARGE Linux} (if not pre-installed)
\begin{itemize}
\item \begin{semiverbatim} apt-get install git \end{semiverbatim}
\item \begin{semiverbatim} yum install git-core \end{semiverbatim}
\end{itemize}

\vspace{2mm}
{\LARGE Mac}
\begin{itemize}
\item With MacPorts: \begin{semiverbatim} sudo port install git-core +svn +doc +bash{\_}completion \end{semiverbatim}
\item With a GUI installer: \url{http://code.google.com/p/git-osx-installer}
\end{itemize}

\vspace{2mm}
{\LARGE Windows}
\begin{itemize}
\item Git bash + GUI installer: \url{http://msysgit.github.com/}
\end{itemize}

\end{frame}

\begin{frame}{Basic Git commands}

\begin{center}
{\huge Demo!}
\end{center}

\end{frame}

\begin{frame}{Git web hosts}

\includegraphics[scale=.03]{octocat} {\LARGE GitHub}
\begin{itemize}
\item Free public repositories
\item Free plan with (limited) private repositories for students and teachers
\item Used VERY widely
\end{itemize}

\vspace{5mm}
\includegraphics[scale=.03]{bitbucketlogo} \hspace{1mm}{\LARGE Bitbucket}
\begin{itemize}
\item Unlimited private repositories for up to 5 users for anyone
\item Unlimited private repositories for {\bf unlimited} users with .edu email address
\item It's great!
\end{itemize}

\vspace{3mm}
(And others)
\end{frame}

\end{document}